\documentclass[a4paper,11pt]{article}
\usepackage[utf8]{inputenc}
\usepackage[T1]{fontenc}
\usepackage[ngerman]{babel}
\usepackage{amsmath,amssymb}
\usepackage{xcolor}
\usepackage{colortbl}
\usepackage{array}
\usepackage{tikz}
\usepackage{enumitem}
\usepackage[margin=2cm]{geometry}
\definecolor{bgdark}{HTML}{1E1E1E}
\definecolor{fgtext}{HTML}{E6E6E6}
\definecolor{gridcol}{HTML}{4A4A4A}
\definecolor{accent}{HTML}{6CB6FF}
\pagecolor{bgdark}
\color{fgtext}
\arrayrulecolor{fgtext}
\setlength{\parindent}{0pt}
\newcommand{\karo}[1]{\par\noindent\begin{tikzpicture}\draw[step=5mm,gridcol,very thin](0,0) grid (\linewidth,#1);\end{tikzpicture}\par\vspace{2mm}}
\newcommand{\aufgabe}[2]{\vspace{3mm}\par\noindent{\large\color{accent}\textbf{Aufgabe #1}}\hfill{\color{gridcol}\normalsize (#2 Punkte)}\par\vspace{1mm}}
\newcommand{\teil}[1]{\vspace{2mm}\par\noindent{\Large\color{accent}\textbf{#1}}\par{\color{gridcol}\rule{\linewidth}{0.8pt}}\par\vspace{1mm}}

\begin{document}
\begin{center}
  {\LARGE\color{accent}\textbf{Optimierungsverfahren}}\\[3pt]
  {\large Probeklausur 05 (XL) --- \textbf{Vollabdeckung}: alle Rechenaufgaben}\\[2pt]
  {\small Deckt den \emph{kompletten} klausurrelevanten Stoff in einer Klausur ab ---
  Grundlagen, Lineare \& Transport, Nichtlinear, Gradientenfrei, Metaheuristiken, Mehrziel.}
\end{center}
\vspace{2mm}
\noindent\textbf{Name:}\ \underline{\hspace{6cm}}\hfill\textbf{Datum:}\ \underline{\hspace{4cm}}\\[4pt]
\noindent\textbf{Gesamtpunkte:}\ 174 \hfill \textbf{Bearbeitungszeit:}\ ca.\ 180 Minuten (XL) \hfill \textbf{Hilfsmittel:}\ Open-Book, Taschenrechner\\[3pt]
{\color{gridcol}\rule{\linewidth}{0.4pt}}
\par\vspace{1mm}
\noindent{\small\textbf{Hinweis:} Übungs-Vollabdeckungsklausur. Die echte Klausur umfasst 50\,P in 60\,min;
hier wird bewusst \emph{jeder} Aufgabentyp einmal geprüft, damit keine Lücke offen bleibt. Begründen Sie kurz.}
\par\vspace{2mm}
\noindent\renewcommand{\arraystretch}{1.2}
\begin{center}\footnotesize
\begin{tabular}{|c|*{14}{c|}}
\hline
Aufg. & 1 & 2 & 3 & 4 & 5 & 6 & 7 & 8 & 9 & 10 & 11 & 12 & 13 & 14\\\hline
Pkt. & 6 & 6 & 5 & 7 & 6 & 6 & 8 & 5 & 20 & 5 & 7 & 8 & 5 & 6\\\hline
\end{tabular}\\[2pt]
\begin{tabular}{|c|*{14}{c|}}
\hline
Aufg. & 15 & 16 & 17 & 18 & 19 & 20 & 21 & 22 & 23 & 24 & 25 & 26 & 27 & 28\\\hline
Pkt. & 6 & 7 & 4 & 5 & 3 & 5 & 4 & 5 & 5 & 6 & 7 & 5 & 5 & 7\\\hline
\end{tabular}\\[2pt]
{\footnotesize Summe $=174$ Punkte}
\end{center}

% ===================== TEIL A: GRUNDLAGEN =====================
\clearpage
\teil{Teil A --- Grundlagen}
\aufgabe{1 --- Modellierung}{6}
Eine Werkstatt fertigt zwei Produkte: Standardstühle ($x_1$) und Designstühle ($x_2$).
Pro Standardstuhl beträgt der Gewinn 40\,Euro, pro Designstuhl 120\,Euro. Täglich stehen
100\,kg Holz und 90\,Arbeitsstunden zur Verfügung. Ein Standardstuhl benötigt 2\,kg Holz und 3\,h,
ein Designstuhl 4\,kg Holz und 2\,h.
\begin{enumerate}[label=(\alph*),nosep,leftmargin=1.8em]
  \item (4\,P) Definieren Sie die Entscheidungsvariablen und stellen Sie das vollständige
        Optimierungsproblem auf (Zielfunktion, Nebenbedingungen, Nichtnegativität).
  \item (2\,P) Handelt es sich um ein lineares Optimierungsproblem? Begründen Sie.
\end{enumerate}
\karo{13cm}

\clearpage
\aufgabe{2 --- Konvexität (Funktion und Menge)}{6}
\begin{enumerate}[label=(\alph*),nosep,leftmargin=1.8em]
  \item (2\,P) Ist $f(x)=x^4$ konvex? Begründen Sie über die zweite Ableitung.
  \item (2\,P) Ist $f(x_1,x_2)=x_1^2+3x_2^2$ konvex? Begründen Sie über die Hesse-Matrix.
  \item (2\,P) Ist die Menge $M=\{\vec x\in\mathbb R^2 : x_1^2+x_2^2\le 4\}$ konvex? Und die Menge
        $N=\{\vec x : x_1^2+x_2^2\ge 4\}$? Begründen Sie kurz.
\end{enumerate}
\karo{13cm}

\clearpage
\aufgabe{3 --- Analytische Optima (eine Variable)}{5}
Gegeben sei $f(x)=x^3-6x^2+9x+1$. Bestimmen Sie alle stationären Punkte, klassifizieren Sie
diese mit dem $f''$-Test als Minimum/Maximum und geben Sie die zugehörigen Funktionswerte an.
\karo{14cm}

\clearpage
\aufgabe{4 --- Analytische Optima (zwei Variablen, Hesse)}{7}
Gegeben sei $f(x_1,x_2)=x_1^2+x_1x_2+x_2^2-3x_1$.
\begin{enumerate}[label=(\alph*),nosep,leftmargin=1.8em]
  \item (3\,P) Bestimmen Sie den stationären Punkt über $\nabla f=\vec 0$ (Gleichungssystem lösen).
  \item (3\,P) Stellen Sie die Hesse-Matrix auf und klassifizieren Sie den Punkt über die Definitheit.
  \item (1\,P) Geben Sie den optimalen Funktionswert an.
\end{enumerate}
\karo{12cm}

% ===================== TEIL B: LINEARE OPTIMIERUNG =====================
\clearpage
\teil{Teil B --- Lineare Optimierung}
\aufgabe{5 --- Grafische Lösung \& Slackvariablen}{6}
\[
\max_{\vec x}\ f(\vec x)=3x_1+2x_2 \quad\text{u.d.N.}\quad
\underbrace{x_1+x_2\le 4}_{g_1},\ \ \underbrace{x_1+3x_2\le 6}_{g_2},\ \ x_1,x_2\ge 0.
\]
\begin{center}
\begin{tikzpicture}[scale=0.85]
  \fill[white,opacity=0.06] (0,0)--(4,0)--(3,1)--(0,2)--cycle;
  \draw[gridcol,very thin,step=1] (0,0) grid (6,5);
  \draw[->,fgtext] (0,0)--(6.4,0) node[right]{$x_1$};
  \draw[->,fgtext] (0,0)--(0,5.4) node[above]{$x_2$};
  \foreach \i in {1,2,3,4,5,6}{\node[fgtext,below] at (\i,0){\small\i};}
  \foreach \j in {1,2,3,4,5}{\node[fgtext,left] at (0,\j){\small\j};}
  \draw[accent,very thick] (4,0)--(0,4) node[above right,pos=0.12]{\small$g_1{:}\,x_1{+}x_2{=}4$};
  \draw[accent,very thick] (6,0)--(0,2) node[below left,pos=0.35]{\small$g_2{:}\,x_1{+}3x_2{=}6$};
\end{tikzpicture}
\end{center}
\begin{enumerate}[label=(\alph*),nosep,leftmargin=1.8em]
  \item (4\,P) Bestimmen Sie die optimale Ecke $\vec x^{\,*}$ und $f(\vec x^{\,*})$.
  \item (2\,P) Welche Werte nehmen die Slackvariablen $s_1,s_2$ im Optimum an? Welche NB ist aktiv?
\end{enumerate}
\karo{8cm}

\clearpage
\aufgabe{6 --- Standardform der Linearen Optimierung}{6}
Überführen Sie in die \textbf{Standardform} (nur Minimierung, nur Gleichungs-NB, alle Variablen
$\ge 0$, rechte Seiten $\ge 0$):
\[
\max_{\vec x}\ 2x_1+3x_2 \quad\text{u.d.N.}\quad
x_1+x_2\le 5,\ \ x_1-x_2\ge 1,\ \ x_1\ \text{frei},\ \ x_2\ge 0.
\]
Geben Sie alle Umformungen explizit an (Slack-/Surplus-Variablen, Aufspaltung freier Variablen, $\max\!\to\!\min$).
\karo{13cm}

\clearpage
\aufgabe{7 --- Simplex (Tableau-Verfahren)}{8}
Lösen Sie mit dem \textbf{Simplex-Tableau} (Schlupfvariablen $s_1,s_2$):
\[
\max_{\vec x}\ f=3x_1+2x_2 \quad\text{u.d.N.}\quad
x_1+x_2\le 4,\ \ x_1+3x_2\le 6,\ \ x_1,x_2\ge 0.
\]
Geben Sie pro Iteration Pivotspalte, Quotiententest, Pivotzeile und das aktualisierte Tableau an.
Nennen Sie $\vec x^{\,*}$ und $f^{*}$.
\karo{13cm}

\clearpage
\aufgabe{8 --- Ganzzahlige Optimierung}{5}
\[
\max_{\vec x}\ f=x_1+x_2 \quad\text{u.d.N.}\quad
2x_1+x_2\le 8,\ \ x_1+2x_2\le 8,\ \ x_1,x_2\ge 0,\ \ x_1,x_2\in\mathbb Z.
\]
\begin{enumerate}[label=(\alph*),nosep,leftmargin=1.8em]
  \item (2\,P) Bestimmen Sie die kontinuierliche optimale Lösung und $f$.
  \item (3\,P) Bestimmen Sie die ganzzahlige optimale Lösung. Liefert Abrunden das Optimum?
\end{enumerate}
\karo{11cm}

% ===================== TEIL C: TRANSPORT =====================
\clearpage
\teil{Teil C --- Transportprobleme}
\aufgabe{9 --- Transportproblem (NWER, MMM, VAM, MODI)}{20}
Drei Werke liefern an drei Läden. Angebote $a=(20,30,25)$, Nachfragen $b=(15,30,30)$
(balanciert, $\sum a=\sum b=75$). Kosten $c_{ij}$:
\begin{center}\renewcommand{\arraystretch}{1.3}
\begin{tabular}{c|ccc|c}
$c_{ij}$ & $D_1$ & $D_2$ & $D_3$ & Angebot\\\hline
$S_1$ & 3 & 5 & 7 & 20\\
$S_2$ & 6 & 4 & 3 & 30\\
$S_3$ & 4 & 2 & 5 & 25\\\hline
Nachfrage & 15 & 30 & 30 &\\
\end{tabular}
\end{center}
\begin{enumerate}[label=(\alph*),nosep,leftmargin=1.8em]
  \item (5\,P) Startlösung mit der \textbf{Nordwesteckenregel} (NWER) + Gesamtkosten.
  \item (5\,P) Startlösung mit der \textbf{Matrix-Minimum-Methode} (MMM) + Gesamtkosten.
  \item (5\,P) Startlösung mit der \textbf{Vogelschen Approximationsmethode} (VAM) + Gesamtkosten.
  \item (5\,P) Prüfen Sie die \textbf{NWER}-Lösung aus (a) mit \textbf{MODI} (Potentiale $u_i,v_j$,
        reduzierte Kosten) auf Optimalität. Falls nicht optimal: eintretende Zelle,
        Stepping-Stone-Zyklus und verbesserte Gesamtkosten nach einem Schritt.
\end{enumerate}
\karo{13cm}
\clearpage
\karo{16cm}

% ===================== TEIL D: NICHTLINEAR =====================
\clearpage
\teil{Teil D --- Nichtlineare Optimierung}
\aufgabe{10 --- Gradientenabstieg}{5}
$f(x_1,x_2)=x_1^2+2x_2^2$, Startpunkt $\vec x_0=(1,1)$, Lernrate $\eta=0{,}1$.
Führen Sie \textbf{zwei} Schritte $\vec x_{k+1}=\vec x_k-\eta\nabla f(\vec x_k)$ durch und geben Sie
jeweils $f$ an.
\karo{12cm}

\clearpage
\aufgabe{11 --- Newton-Verfahren (1D und 2D)}{7}
\begin{enumerate}[label=(\alph*),nosep,leftmargin=1.8em]
  \item (3\,P) \textbf{1D:} $f(x)=x^3-6x^2+9x$, ein Newton-Schritt $x_{k+1}=x_k-\dfrac{f'(x_k)}{f''(x_k)}$
        mit $x_0=4$.
  \item (4\,P) \textbf{2D:} $f(x_1,x_2)=x_1^2+x_1x_2+x_2^2-3x_1$, ein Newton-Schritt
        $\vec x_{k+1}=\vec x_k-H^{-1}\nabla f(\vec x_k)$ mit $\vec x_0=(0,0)$ (inkl.\
        $2\times2$-Inverse). Was fällt auf?
\end{enumerate}
\karo{13cm}

\clearpage
\aufgabe{12 --- KKT-Bedingungen aufstellen und lösen}{8}
\[
\min_{\vec x}\ f(\vec x)=(x_1-2)^2+(x_2-2)^2 \quad\text{u.d.N.}\quad g(\vec x)=x_1+x_2-2\le 0 .
\]
\begin{enumerate}[label=(\alph*),nosep,leftmargin=1.8em]
  \item (4\,P) Lagrange-Funktion und alle vier KKT-Bedingungen aufstellen.
  \item (4\,P) Lösen mit Fallunterscheidung ($\mu=0$ vs.\ $g=0$), Vorzeichen prüfen, $\vec x^{\,*}$ und $f^{*}$ angeben.
\end{enumerate}
\karo{13cm}

\clearpage
\aufgabe{13 --- Grafische nichtlineare Optimierung}{5}
Minimiert werden soll $f(\vec x)=(x_1-3)^2+(x_2-3)^2$ (kreisförmige Isolinien um $(3,3)$, blau
gestrichelt) im zulässigen Bereich $x_1+x_2\le 4,\ x_1,x_2\ge 0$ (grau). Bestimmen Sie grafisch
$\vec x^{\,*}$ und $f(\vec x^{\,*})$.
\begin{center}
\begin{tikzpicture}[scale=0.85]
  \fill[white,opacity=0.06] (0,0)--(4,0)--(0,4)--cycle;
  \draw[gridcol,very thin,step=1] (0,0) grid (5,5);
  \draw[->,fgtext] (0,0)--(5.4,0) node[right]{$x_1$};
  \draw[->,fgtext] (0,0)--(0,5.4) node[above]{$x_2$};
  \foreach \i in {1,2,3,4,5}{\node[fgtext,below] at (\i,0){\small\i}; \node[fgtext,left] at (0,\i){\small\i};}
  \draw[accent,very thick] (4,0)--(0,4) node[above right,pos=0.1]{\small$x_1{+}x_2{=}4$};
  \fill[red!80!white] (3,3) circle(2.4pt) node[above right,fgtext]{\small$(3,3)$};
  \draw[blue!60!white,dashed] (3,3) circle(1.0);
  \draw[blue!60!white,dashed] (3,3) circle(1.414);
  \draw[blue!60!white,dashed] (3,3) circle(2.0);
\end{tikzpicture}
\end{center}
\karo{8cm}

% ===================== TEIL E: GRADIENTENFREI =====================
\clearpage
\teil{Teil E --- Gradientenfreie Verfahren}
\aufgabe{14 --- Nelder-Mead}{6}
Für $f(x_1,x_2)=x_1^2+x_2^2$ ist der sortierte Simplex gegeben ($x_0$ bester, $x_2$ schlechtester),
$\alpha=1$:
\[
x_0=(1,1),\ f=2;\qquad x_1=(3,1),\ f=10;\qquad x_2=(2,3),\ f=13.
\]
Berechnen Sie Schwerpunkt $x_c$ (ohne $x_2$) und Reflexionspunkt $x_r$. Bestimmen Sie über $f(x_r)$,
welcher Schritt als nächstes folgt.
\karo{12cm}

\clearpage
\aufgabe{15 --- DIRECT: potentiell optimale Rechtecke}{6}
DIRECT-Durchlauf (Minimierung), fünf Rechtecke mit Größe $d_i$ und Mittelpunktswert $f(c_i)$.
Bestimmen Sie die \textbf{potentiell optimalen} Rechtecke (untere-rechte konvexe Hülle) und begründen Sie.
\begin{center}\renewcommand{\arraystretch}{1.2}
\begin{tabular}{c|ccccc}
Nr.\ $i$ & 1 & 2 & 3 & 4 & 5\\\hline
$d_i$ & 1 & 1 & 3 & 3 & 6\\
$f(c_i)$ & 6 & 9 & 4 & 7 & 2\\
\end{tabular}
\end{center}
\begin{center}
\begin{tikzpicture}[scale=0.7]
  \draw[gridcol,very thin,step=1] (0,0) grid (7,10);
  \draw[->,fgtext] (0,0)--(7.4,0) node[right]{$d$};
  \draw[->,fgtext] (0,0)--(0,10.4) node[above]{$f(c)$};
  \foreach \i in {1,2,3,4,5,6}{\node[fgtext,below] at (\i,0){\small\i};}
  \foreach \j in {2,4,6,8,10}{\node[fgtext,left] at (0,\j){\small\j};}
  \foreach \n/\x/\y in {1/1/6,2/1/9,3/3/4,4/3/7,5/6/2}{
    \fill[accent](\x,\y) circle(2.6pt); \node[fgtext,above right] at (\x,\y){\small \n};}
\end{tikzpicture}
\end{center}
\karo{7cm}

\clearpage
\aufgabe{16 --- Quadratische (QO) und Sequentielle Quadratische Optimierung (SQO)}{7}
\begin{enumerate}[label=(\alph*),nosep,leftmargin=1.8em]
  \item (4\,P) Schreiben Sie $f(x_1,x_2)=x_1^2+x_1x_2+2x_2^2-x_1$ als
        $\tfrac12\vec x^{\!\top}Q\vec x+\vec c^{\!\top}\vec x$. Geben Sie $Q,\vec c$ an, prüfen Sie
        $Q$ auf positive Definitheit und bestimmen Sie das Minimum.
  \item (3\,P) Erläutern Sie das Grundprinzip der SQO: Was wird je Iteration gelöst, und wie hängt
        es mit der QO zusammen?
\end{enumerate}
\karo{12cm}

\clearpage
\aufgabe{17 --- SMBO / Surrogatmodell}{4}
Eine sehr teure Zielfunktion soll minimiert werden. Beschreiben Sie die vier Schritte der
SMBO-Schleife. Warum ein Surrogat, und was geschieht damit nach jeder echten Auswertung?
\karo{14cm}

% ===================== TEIL F: METAHEURISTIKEN =====================
\clearpage
\teil{Teil F --- Metaheuristiken}
\aufgabe{18 --- $(1+1)$-Evolutionsstrategie (1/5-Erfolgsregel)}{5}
Minimierung. Speicherlänge $g=10$, Multiplikator $a=1{,}2$, $\sigma=0{,}5$. Aktuelle Qualität
$f(x)=0{,}50$, mutierte $f(x^{\text{next}})=0{,}48$. Inkl.\ dieses Schritts liegen 3 Erfolge im Fenster.
\begin{enumerate}[label=(\alph*),nosep,leftmargin=1.8em]
  \item (2\,P) War der Schritt ein Erfolg? Wird $x$ übernommen?
  \item (3\,P) Berechnen Sie die Erfolgsrate $s$ und die neue Schrittweite $\sigma$.
\end{enumerate}
\karo{11cm}

\clearpage
\aufgabe{19 --- CMA-ES (Konzept)}{3}
Was unterscheidet die CMA-ES von der einfachen $(1+1)$-ES? Was wird fortlaufend angepasst und
welchen Vorteil bringt das?
\karo{14cm}

\clearpage
\aufgabe{20 --- Simulated Annealing}{5}
\begin{enumerate}[label=(\alph*),nosep,leftmargin=1.8em]
  \item (3\,P) Minimierung: Nachbar um $\Delta f=3$ schlechter, $T=5$. Akzeptanzwahrscheinlichkeit
        $p=\exp(-\Delta f/T)$? Wird sie bei $u=0{,}5$ angenommen?
  \item (2\,P) Rolle der Temperatur $T$ zu Beginn und am Ende?
\end{enumerate}
\karo{12cm}

\clearpage
\aufgabe{21 --- Particle Swarm Optimization (PSO)}{4}
Geschwindigkeits-Update $v\leftarrow w\,v+c_1 r_1(p-x)+c_2 r_2(g-x)$, dann $x\leftarrow x+v$.
Gegeben (1D): $x=2$, $v=1$, $p=0$, $g=-1$, $w=0{,}5$, $c_1=c_2=1$, $r_1=r_2=1$.
\begin{enumerate}[label=(\alph*),nosep,leftmargin=1.8em]
  \item (1\,P) Benennen Sie die drei Komponenten (Trägheit, kognitiv, sozial).
  \item (3\,P) Berechnen Sie neue Geschwindigkeit $v$ und Position $x$.
\end{enumerate}
\karo{11cm}

\clearpage
\aufgabe{22 --- Ant Colony Optimization (ACO)}{5}
Auswahl zwischen Kanten $A,B$ mit $p_e\propto \tau_e^{\alpha}\eta_e^{\beta}$. Gegeben
$\tau_A=2,\eta_A=1$ und $\tau_B=1,\eta_B=4$, $\alpha=\beta=1$.
\begin{enumerate}[label=(\alph*),nosep,leftmargin=1.8em]
  \item (3\,P) Berechnen Sie $p_A,p_B$.
  \item (2\,P) Aktualisieren Sie $\tau_A$ mit Verdunstung $\rho=0{,}5$ und Zuwachs $\Delta\tau_A=1$
        gemäß $\tau\leftarrow(1-\rho)\tau+\Delta\tau$.
\end{enumerate}
\karo{11cm}

\clearpage
\aufgabe{23 --- Evolutionäre Operatoren (binär)}{5}
Eltern $x_1=(1,1,0,0,1,1,0,0)$, $x_2=(0,0,1,1,0,0,1,1)$.
\begin{enumerate}[label=(\alph*),nosep,leftmargin=1.8em]
  \item (1\,P) Bitflip-Mutation von $x_1$ an Position 3.
  \item (2\,P) 1-Punkt-Crossover (Schnitt nach Position 4).
  \item (2\,P) 2-Punkt-Crossover (Schnitte nach Position 2 und 6).
\end{enumerate}
\karo{12cm}

% ===================== TEIL G: MEHRZIEL =====================
\clearpage
\teil{Teil G --- Mehrzieloptimierung}
\aufgabe{24 --- Dominanz \& Pareto-Menge}{6}
Punkte im Zielraum (beide Ziele \emph{minimieren}):
$A=(2,6),\ B=(3,4),\ C=(5,3),\ D=(4,7),\ E=(6,2)$.
\begin{enumerate}[label=(\alph*),nosep,leftmargin=1.8em]
  \item (2\,P) Geben Sie alle Dominanzbeziehungen an.
  \item (4\,P) Bestimmen Sie die Pareto-optimale Menge (Pareto-Front).
\end{enumerate}
\karo{12cm}

\clearpage
\aufgabe{25 --- Non-Dominated-Sorting-Rang}{7}
Punkte (beide Ziele \emph{minimieren}):
$P_1=(1,4),\ P_2=(2,2),\ P_3=(4,1),\ P_4=(2,5),\ P_5=(3,3),\ P_6=(5,4)$.
Bestimmen Sie für jeden Punkt den Non-Dominated-Sorting-Rang (schichtweise).
\karo{13cm}

\clearpage
\aufgabe{26 --- Crowding Distance}{5}
Front (beide Ziele \emph{minimieren}): $A=(1,5),\ B=(2,4),\ C=(3,2),\ D=(5,1)$.
\begin{enumerate}[label=(\alph*),nosep,leftmargin=1.8em]
  \item (4\,P) Berechnen Sie die Crowding Distance jedes Punktes (Normierung über $f_{\max}-f_{\min}$
        pro Ziel, Randpunkte $=\infty$).
  \item (1\,P) Welche zwei Punkte würde ein EA zuerst auswählen?
\end{enumerate}
\karo{12cm}

\clearpage
\aufgabe{27 --- NSGA-II vs.\ SMS-EMOA}{5}
Beschreiben Sie den Ablauf des NSGA-II (Selektion per Rang, bei Gleichstand Crowding Distance)
und den zentralen Unterschied zum SMS-EMOA (Sekundärkriterium).
\karo{14cm}

\clearpage
\aufgabe{28 --- Hypervolumen \& Hypervolumen-Beitrag}{7}
Front (beide Ziele \emph{minimieren}): $P_1=(1,4),\ P_2=(2,2),\ P_3=(4,1)$, Referenzpunkt $r=(5,5)$.
\begin{enumerate}[label=(\alph*),nosep,leftmargin=1.8em]
  \item (4\,P) Berechnen Sie das Hypervolumen.
  \item (3\,P) Berechnen Sie den Hypervolumen-Beitrag (HVC) jedes Punktes.
\end{enumerate}
\karo{12cm}

\vfill
\begin{center}{\color{gridcol}\small--- Ende Probeklausur 05 (Vollabdeckung, 174 Punkte) ---}\end{center}
\end{document}
